\documentclass[11pt]{article}
\usepackage[margin=1in]{geometry}
\usepackage{amsmath}
\usepackage{booktabs}
\usepackage{hyperref}

\newcommand{\col}[1]{\texttt{\detokenize{#1}}}

\title{Brightness and noise nuisances in the age-correlation analysis}
\author{}
\date{October 2026}

\begin{document}
\maketitle

\section{Summary}

\begin{itemize}
  \item The nuisance \col{log_med_err_frac} (median fractional flux error) is, to within
        a few per cent, a measure of aperture brightness: it scales as $F^{-1/2}$, as photon
        noise should.
  \item In the LMC aperture brightness is itself a strong age correlator (younger clusters
        have brighter apertures, largely through the mass-to-light ratio). So
        \col{log_med_err_frac} carries almost all of the nuisance-only age prediction in the
        LMC, and conditioning on it removes real (but non-variability) age information.
  \item The light-curve features that describe noise structure (red-noise normalisation and
        slope, GP jitter, white-noise ratio, Stetson~$J$, von~Neumann ratio) correlate with
        \col{log_med_err_frac} about as strongly as with age.
  \item We reparameterise the noise term into a brightness term \col{log_F} and a noise term
        that is (nearly) independent of brightness, \col{log_noise_perp}, so brightness can
        be switched on and off. ``On'' measures age information beyond anything brightness
        related (variability specific); ``off'' measures all physical correlators.
  \item Separately: \col{log_med_flux} is identically zero (the light curves are normalised
        before it is computed) and carries no information.
\end{itemize}

Unless stated otherwise, numbers are for one row per cluster (the median over all of a
cluster's sector combinations; ``setup A'' in \col{08_simple_overview.ipynb}), LMC $n=118$,
MW $n=124$, SMC $n=106$. Ridge $R^2$ values are 5-fold cross-validated, averaged over 20
random fold assignments, with \texttt{StandardScaler} + \texttt{RidgeCV} as in the notebook.

\section{What the original nuisances measure}

\paragraph{\col{log_med_flux} is dead.}
The pipeline divides every sector by its own median flux before storing the light curve
(\col{normalize_flux} in \col{lc_lsp.py}), so the median of the stored flux is 1 and its
logarithm is 0 for all 17\,233 rows. It never acted as a brightness control.

\paragraph{\col{log_med_err_frac} is brightness.}
It is the median of the normalised errors, $\sigma/F$, of the (30-min resampled) light curve.
Normalising does not remove the brightness dependence; it creates it. For photon noise
$\sigma \propto \sqrt{F}$, so
\begin{equation}
  \frac{\sigma}{F} \propto F^{-1/2}.
\end{equation}
Measured against the median raw ELK aperture flux $F$:
\begin{center}
\begin{tabular}{lccc}
\toprule
 & MW & LMC & SMC \\
\midrule
slope $d\log(\sigma/F)/d\log F$ & $-0.524$ & $-0.512$ & $-0.496$ \\
Spearman $\rho(\sigma/F,\ F)$ & $-0.980$ & $-0.997$ & $-0.997$ \\
\bottomrule
\end{tabular}
\end{center}

\section{Why this matters: brightness is an age correlator}

In the LMC (one distance) the aperture flux is $F = F_{\rm cl} + F_{\rm field} + F_{\rm sky}$, and
the cluster light at fixed mass falls with age as the TESS-band mass-to-light ratio
$\Upsilon_T$ rises. Hence $\rho(\log F, \text{age}) = -0.72$ and
$\rho(\log\sigma/F, \text{age}) = +0.71$ in the LMC. (In the MW, brightness is driven mainly
by distance and field and barely tracks age, $\rho = +0.13$.)

LMC ridge regression of log age on nuisances only:
\begin{center}
\begin{tabular}{lc}
\toprule
nuisance set & CV $R^2$ \\
\midrule
all original nuisances & 0.45 \\
original without \col{log_med_err_frac} & 0.11 \\
\col{log_med_err_frac} alone & 0.48 \\
\bottomrule
\end{tabular}
\end{center}
So essentially all of the nuisance-only age prediction in the LMC is ``younger clusters
have brighter apertures''. This is real astrophysical signal (see also the cluster light
fraction vs age analysis in \col{GenerateTable_CGW/crowding_analysis/}), but it is about
brightness, not about variability.

\paragraph{Two routes through brightness.}
Brightness reaches the features in two ways:
\begin{enumerate}
  \item \emph{astrophysical}: age $\to$ cluster light $\to F$ (a real age correlator);
  \item \emph{measurement}: $F \to$ photon-noise floor $\sigma/F \to$ measured noise-structure
        features (detectability, continuum level, jitter).
\end{enumerate}
Both pass through the same variable $F$, so no choice of nuisances can separate them.
Controlling for brightness removes both; not controlling keeps both.

\section{The features already contain the brightness information}

LMC, Spearman correlations (selected features; full table from the script in
Section~\ref{sec:impl}). ``Partial'' columns residualise both feature and age on the
nuisance set by OLS. ``base'' is the six nuisances with no brightness or noise term
(cadence, $n_{\rm sectors}$, $n_{\rm points}$, radius, Dec, $b$).
\begin{center}
\begin{tabular}{lrrrrr}
\toprule
feature & $\rho(\sigma/F)$ & raw $\rho$(age) & base & brightness off & brightness on \\
\midrule
\col{gp_jitter}            & $-0.63$ & $-0.60$ & $-0.50$ & $-0.51$ & $-0.29$ \\
\col{rn_log10_N}           & $-0.62$ & $-0.58$ & $-0.45$ & $-0.44$ & $-0.20$ \\
\col{rn_alpha}             & $-0.58$ & $-0.59$ & $-0.47$ & $-0.47$ & $-0.24$ \\
\col{stetson_j_gap_aware}  & $-0.54$ & $-0.49$ & $-0.35$ & $-0.34$ & $-0.16$ \\
\col{gp_rn_wn_ratio}       & $-0.53$ & $-0.52$ & $-0.44$ & $-0.43$ & $-0.26$ \\
\col{vn_ratio_gap_aware}   & $-0.48$ & $-0.47$ & $-0.26$ & $-0.25$ & $-0.15$ \\
\col{gamma_p_sec}          & $+0.22$ & $+0.26$ & $+0.27$ & $+0.27$ & $-0.00$ \\
\col{mean_median_offset}   & $+0.22$ & $+0.26$ & $+0.22$ & $+0.22$ & $-0.00$ \\
\col{intrinsic_std}        & $-0.20$ & $-0.25$ & $-0.24$ & $-0.23$ & $-0.17$ \\
\bottomrule
\end{tabular}
\end{center}
\begin{itemize}
  \item Noise-structure features roughly halve when brightness is controlled, but stay
        non-zero ($-0.15$ to $-0.29$).
  \item Skewness/offset features vanish: their age trend is entirely brightness mediated.
  \item Amplitude features (\col{intrinsic_std}, \col{excess_var}) barely depend on
        brightness and move least ($-0.23 \to -0.17$).
\end{itemize}

LMC ridge, nuisances (N) vs nuisances + features (N+F):
\begin{center}
\begin{tabular}{lccc}
\toprule
nuisance set & N & N+F & gain \\
\midrule
original           & 0.45 & 0.47 & $+0.02$ \\
brightness on      & 0.45 & 0.47 & $+0.02$ \\
brightness off     & 0.10 & 0.46 & $+0.36$ \\
base               & 0.11 & 0.46 & $+0.35$ \\
features only      &      & 0.50 &  \\
\bottomrule
\end{tabular}
\end{center}
The total age information is the same ($R^2 \approx 0.46$--$0.50$) whichever nuisances are
used; only its attribution moves. Features alone predict age as well as brightness alone
(0.50 vs 0.48), and combining them adds nothing. In the MW and SMC all of these $R^2$ are
$\lesssim 0.1$ and removing \col{log_med_err_frac} changes the partial correlations by at most
$\sim 0.1$.

\section{Reparameterisation}

All three quantities are computed per sector combination (\col{add_brightness_nuisances} in
\col{analysis_utils.py}).
\begin{description}
  \item[\col{log_F}] $\log_{10}$ of the median, over the combination's sectors, of each
        sector's median raw (un-normalised) ELK flux [e/s]. Per cluster it reproduces
        \col{F_meas} of the crowding analysis exactly.
  \item[\col{log_noise_q}] $\log_{10}(\sigma/F) + \tfrac12\log_{10}F = \log_{10}(\sigma/\sqrt{F})$:
        the noise in units of the photon noise expected for that flux. It uses the error of
        the light curve actually analysed (after resampling), so it does not pick up the
        cadence dependence of the raw, unbinned errors.
  \item[\col{log_noise_perp}] $\log_{10}(\sigma/F) - (a + b\log_{10}F)$, with $a,b$ an OLS fit per
        galaxy on the per-cluster medians, applied to every row. Fitted slopes:
        $b = -0.524$ (MW), $-0.512$ (LMC), $-0.496$ (SMC). Linearly uncorrelated with
        $\log F$ by construction. Coefficients can be fixed with \col{perp_slopes}.
\end{description}
For a linear model, $\{\log F, \text{noise term}\}$ spans the same space as
$\{\log F, \log\sigma/F\}$, which is (to $\rho = -0.997$) the space of $\log\sigma/F$ alone.
So \emph{brightness on} reproduces the original analysis (ridge LMC $R^2$ 0.45 vs 0.45; feature
partial correlations differ by at most 0.016), while \emph{brightness off} drops the
brightness direction only. The reparameterisation is a relabelling; its purpose is to make
brightness a switch. Tree models (random forest) are not invariant to such recombinations,
so for them ``on'' will be close to, but not identical to, ``original''.

\section{How clean is the decorrelation?}

\begin{center}
\begin{tabular}{llccc}
\toprule
 & & MW & LMC & SMC \\
\midrule
\col{log_noise_q}    & $\rho(\cdot, \log F)$ & $-0.37$ & $-0.46$ & $+0.07$ \\
                     & $\rho(\cdot, \text{age})$ & $-0.06$ & $+0.25$ & $+0.40$ \\
                     & s.d.\ [dex] & 0.029 & 0.009 & 0.009 \\
\midrule
\col{log_noise_perp} & $\rho(\cdot, \log F)$ & $+0.11$ & $+0.17$ & $-0.05$ \\
                     & $\rho(\cdot, \text{age})$ & $-0.01$ & $-0.26$ & $+0.38$ \\
                     & s.d.\ [dex] & 0.028 & 0.008 & 0.009 \\
\bottomrule
\end{tabular}
\end{center}

\paragraph{Why \col{log_noise_q} still tracks brightness.}
With a sky term, $\sigma^2 \propto F + B$, so $\sigma/\sqrt{F} = \sqrt{1 + B/F}$: faint apertures
are relatively more sky dominated. Any noise ratio defined against source flux alone will
correlate with brightness. The effect is tiny in absolute terms: \col{log_noise_q} spreads by
only $\sim$1--3\%, and the mismatch between the fitted slope ($-0.51$) and $-0.5$ over the
$\sim$2 dex flux range is comparable to that whole spread, which is why a rank correlation
looks large.

\paragraph{\col{log_noise_perp} is only linearly orthogonal.}
Its rank correlation with $\log F$ is $+0.17$ in the LMC. Quadratic and cubic fits do not
reduce it, so it is not simple curvature. A sky-corrected version,
$\sigma/\sqrt{F - F_{\rm sky}}$ with the zodiacal model, was worse ($\rho(\log F) = +0.78$) and is
not used.

\paragraph{Does the leak matter?}
In the LMC, adding either noise term to the base nuisances changes nuisance-only ridge
$R^2$ from 0.11 to 0.10 (\col{log_noise_perp}) or 0.14 (\col{log_noise_q}), and shifts the
feature partial correlations by at most 0.02 (\col{log_noise_perp}) or 0.06
(\col{log_noise_q}). The leak is negligible for the LMC.

\paragraph{SMC caveat.}
In the SMC the noise term correlates with age ($\rho \approx +0.4$) and raises nuisance-only
ridge $R^2$ from $-0.04$ to $0.11$. This comes through the sky fraction
($\rho(\text{noise}, F_{\rm sky}/F) \approx +0.8$): the SMC apertures are twice the intended
radius (Bica diameters used as radii) and collect a lot of sky. So in the SMC the
``brightness off'' specification still contains a small age-linked measurement nuisance.
Report it.

\section{Interpreting the switch}

\begin{description}
  \item[brightness on] (\col{NUIS_CHOICE = 'brightness_on'}): age information beyond anything
        brightness related: the \emph{variability-specific} result.
  \item[brightness off] (\col{'brightness_off'}): all age information, including what brightness
        carries.
\end{description}
``Off minus on'' is \emph{brightness-mediated} information. It is not purely physical: it also
contains the measurement route (photon-noise floor changing the noise-structure features).
To show that the difference is astrophysical, a feature-level test is needed, e.g.\ inject
white noise into bright LMC light curves to match the $\sigma/F$ of faint ones, recompute the
features, and measure how far they move. That isolates the measurement part.

\section{Caveats and next steps}
\begin{itemize}
  \item \textbf{Noise-injection test} (above) to split brightness-mediated information into
        measurement and astrophysical parts.
  \item \textbf{Dilution.} The cluster light fraction $f_{\rm cl}$ is also a brightness-type correlator
        (it carries the mass-to-light ratio) but is not redundant with $F$. It could be a third
        switch: measurement only $\to$ $+$brightness $\to$ $+$brightness$+$dilution. It shares the
        over-control issue.
  \item \textbf{Sky.} $F$ includes the zodiacal light, which carries no cluster information;
        subtracting a reliable sky estimate would make \col{log_F} purely astrophysical.
  \item \textbf{Galactic latitude} enters as signed $b$; $|b|$ is the crowding proxy in the MW.
\end{itemize}

\section{Implementation}\label{sec:impl}
\begin{itemize}
  \item \col{analysis_utils.py}: \col{BRIGHTNESS_SPEC}, \col{sector_raw_flux()} (reads the ELK
        files once and caches per-sector raw medians in \col{sector_raw_flux.csv}), and
        \col{add_brightness_nuisances(rows, perp_slopes=None)}. Not called by \col{load_rows},
        so other notebooks are unchanged; \col{cluster_table} aggregates the new columns when
        present.
  \item \col{08_simple_overview.ipynb}: \col{NUIS_SETS} with \col{'original'},
        \col{'brightness_on'}, \col{'brightness_off'} and the \col{_q} variants; choose with
        \col{NUIS_CHOICE}. \col{load_combos} adds the new columns; \col{per_cluster} keeps every
        candidate nuisance.
\end{itemize}

\end{document}
